\begin{textsample}{POS Dim 2 – 2022\_c4 – Score 22.00 – 2022\_c4/enade\_2022\_tec\_gest\_publ\_f3d64ddc4a.txt}  \label{ex:f2_pos_025}
CAPÍTULO 5 PERCEPÇÃO DA PROVA Nas análises feitas neste capítulo, constam as percepções dos concluintes da área de Tecnologia em Gestão Pública sobre a prova aplicada no Enade/2022.
Essas percepções foram mensuradas por meio de nove questões que avaliaram desde o grau de dificuldade da prova até o tempo gasto para concluí -la.
As percepções sobre a prova foram relacionadas ao desempenho dos estudantes, à grande região de funcionamento do curso, à categoria administrativa e ao tipo de organização acadêmica da IES.
O \textbf{questionário} de percepção da prova encontra -se ao final do Anexo VIII, no qual está a \textbf{reprodução} do exame.
Na apresentação dos dados relativos às nove questões sobre as percepções a respeito da prova, o desempenho dos estudantes foi classificado em quatro quartos.
Para tanto, esse desempenho foi \textbf{ordenado} de forma ascendente.
O percentil 25, P25, também conhecido como primeiro quartil, é a nota de desempenho que deixa um quarto ( 25 \% ) dos valores observados abaixo e três quartos acima.
O quarto inferior de desempenho é composto pelas notas abaixo do primeiro quartil.
Já o percentil 75, P75, também conhecido como terceiro quartil, é o valor para o qual há três quartos ( 75 \% ) dos dados abaixo e um quarto acima dele.
O quarto superior de desempenho é composto pelas notas iguais ou acima do terceiro quartil.
O percentil 50, P50, também conhecido como mediana, é o valor que divide as notas em dois conjuntos de igual tamanho.
O segundo quarto inclui os valores entre o primeiro quartil ( P25 ) e a mediana.
O terceiro quarto contém os valores entre a mediana ( P50 ) e o terceiro quartil ( P75 ).
Vale ressaltar que percentis, quartis e medianas são pontos que, não obrigatoriamente, pertencem ao conjunto original de dados, ao passo que os quartos são \textbf{subconjuntos} dos dados originais.
Na figura 5.1, apresenta -se uma ilustração dos quatro quartos descritos.
A seguir, serão apresentados gráficos com resultados selecionados, relativos às nove questões avaliadas por grupos de estudantes.
Nas barras dos gráficos, apresenta -se o percentual de estudantes que assinalaram uma das opções ou a soma das porcentagens daqueles que assinalaram duas ( ou três ) delas.
Por exemplo, para as questões 1 e 2, nos gráficos, é apresentada a porcentagem total de participantes que assinalaram as opções ( D ) “ Difícil ” e ( E ) “ Muito difícil ”.
Em cada barra, foram assinalados, também, os extremos do intervalo de confiança de 95 \% como linhas verticais unidas por uma linha horizontal na forma da letra H maiúscula.
O \textbf{estimador} de um \textbf{parâmetro} com um certo nível de confiança ( por exemplo, 95,0 \% ) deve conter o \textbf{parâmetro} no intervalo de confiança em 95 \% das vezes.
Na comparação entre os \textbf{estimadores} dos \textbf{parâmetros} de duas classes de uma dada categoria ( por exemplo : região Norte e região Nordeste nas grandes regiões, de primeiro e último quarto dentro de desempenho, de IES públicas e privadas, ou de universidades e faculdades ), associados aos seus respectivos intervalos de confiança, diz -se que não há diferença estatisticamente significativa entre os \textbf{parâmetros} das duas categorias quando há \textbf{interseção} entre os intervalos de confiança, e que há diferença, se os intervalos de confiança são \textbf{disjuntos} ( para maiores informações \textbf{vide} glossário ).
Nas Tabelas do Anexo V, são apresentados os valores absolutos e a distribuição percentual⁹ das alternativas válidas das nove questões, segundo o mesmo \textbf{recorte} de desempenho dos estudantes e da grande região de funcionamento do curso, de categoria administrativa e do tipo de organização acadêmica da IES. 5.1 GRAU DE DIFICULDADE DA PROVA Nesta seção, são apresentadas as percepções dos estudantes quanto ao grau de dificuldade da prova Enade/2022, que foram analisadas considerando sua vinculação ao Componente de Formação Geral e ao Componente Específico. ⁹ \textbf{Reitera} -se que uma das convenções para tabelas \textbf{numéricas} ( Subseção 2.2.3 do Capítulo 2 ) refere -se à possibilidade de, por questão de \textbf{arredondamento}, a soma das partes não resultar em 100 \%. 5.1.1 COMPONENTE DE FORMAÇÃO GERAL O Gráfico 5.1 apresenta as respostas obtidas para a questão 1 : “ Qual o grau de dificuldade desta prova na parte de Formação Geral? ”.
De acordo com os dados do Gráfico 5.1, a região Norte apresentou o maior percentual de estudantes que consideraram a prova “ Difícil ” ou “ Muito difícil ”, com 39,3 \%, seguida da região Sudeste com 38,5 \%.
Na região Nordeste a proporção foi de 35,6 \%, na Região Sul, de 36,8 \%, e na região Centro-Oeste, de 34,3 \%, esta última representando o menor percentual.
O Gráfico 5.2 traz mais resultados sobre a questão tratada no gráfico anterior, porém por quartos de desempenho.
Os dados do Gráfico 5.2 apresentam que o percentual de estudantes que consideraram o componente de Formação Geral da prova “ Difícil ” ou “ Muito difícil ” foi decrescente em função dos quartos de desempenho : 43,5 \%, no 1º quarto inferior e 28,5 \%, no 4º quarto superior, o grupo de melhor desempenho na prova.
Nos quartos de desempenho intermediários, a proporção de estudantes que consideraram a prova “ Difícil ” ou “ Muito difícil ” foi de 38,7 \%, no 2º quarto, e de 36,8 \%, no 3º quarto.
O Gráfico 5.3 traz outros resultados também em relação à mesma questão retratada no gráfico anterior, porém por organização acadêmica das IES.
Como apresentam os dados do Gráfico 5.3, o percentual de estudantes que consideraram o componente de Formação Geral da prova “ Difícil ” ou “ Muito difícil ” foi menor entre os estudantes de CEFET/IF, alcançando 29,4 \%, enquanto, entre os estudantes das faculdades e dos centros universitários, o percentual alcançou 37,9 \% e 36 \%, respectivamente.
O maior percentual foi notado entre os estudantes de universidades 38,1 \%.
O Gráfico 5.4 apresenta os dados por categoria administrativa das IES.
Os dados do Gráfico 5.4 apresentam que o percentual de estudantes que consideraram o Componente de Formação Geral da prova “ Difícil ” ou “ Muito difícil ” foi maior entre os estudantes de cursos de IES privadas ( 37,4 \% ). 5.1.2 COMPONENTE ESPECÍFICO O Gráfico 5.5 apresenta as respostas obtidas para a Questão 2 “ Qual o grau de dificuldade desta prova na parte de Componente Específico? ”.
De acordo com os dados do Gráfico 5.5, a região Norte apresentou o maior percentual de estudantes que consideraram o componente específico da prova “ Difícil ” ou “ Muito difícil ”, com 43,2 \%, seguida da região Sudeste com 41,5 \%.
Na região Nordeste e região Sul, a média foi de 40,1 \% e de 40,9 \%, respectivamente.
A menor média foi constatada na região Centro-Oeste com 39 \%.
O Gráfico 5.6 traz mais resultados sobre a questão tratada no gráfico anterior, porém por quartos de desempenho.
Os dados do Gráfico 5.6 apresentam que o percentual de estudantes que consideraram o componente Específico da prova “ Difícil ” ou “ Muito difícil ” foi de 41,0 \%, no 1.º quarto inferior, e de 41,8 \%, no 3.º quarto superior.
O último quarto superior apresenta ligeira queda no percentual, atingindo 39,0 \%.
O Gráfico 5.7 traz outros resultados também em relação à mesma questão retratada no gráfico anterior, porém por organização acadêmica das IES.
Como apresentam os dados do Gráfico 5.7, o percentual de estudantes que consideraram o Componente Específico da prova “ Difícil ” ou “ Muito difícil ” foi maior entre os estudantes de universidades, alcançando 41,5 \%, enquanto, entre os estudantes dos centros universitários e dos CEFET/IF, as organizações acadêmicas com a menores incidências, alcançaram -se os percentuais de 40,1 \% e 38 \%, respectivamente.
O Gráfico 5.8 apresenta os dados por categoria administrativa das IES.
Os dados do Gráfico 5.8 apresentam que o percentual de estudantes que consideraram o Componente Específico da prova “ Difícil ” ou “ Muito difícil ” foi ligeiramente maior entre os estudantes de cursos de IES privadas ( 41 \% ).
Enquanto a proporção registrada em IES públicas foi de 38,1 \%. 5.2 EXTENSÃO DA PROVA EM RELAÇÃO AO TEMPO TOTAL O Gráfico 5.9 apresenta as respostas obtidas para a questão 3 “ Considerando a extensão da prova, em relação ao tempo total, você considera que a prova foi longa ou muito longa? ”.
Em relação aos dados do Gráfico 5.9, entre as grandes regiões, os que avaliaram a prova como “ Longa ” ou “ Muito longa ”, em relação ao tempo total destinado à sua resolução, atingiram uma porcentagem que variou de 34,7 \%, na região Centro-Oeste, até 39,3 \%, na região Sul, região com maior percentual.
O Gráfico 5.10 traz mais resultados sobre a questão tratada no gráfico anterior, porém por quartos de desempenho.
No Gráfico 5.10, observa -se que a proporção de estudantes que consideraram a prova “ Longa ou “ Muito longa ” foi menor no 4.º quarto de desempenho superior ( 35,4 \% ), e maior no 1.º quarto inferior ( 39,9 \% ).
Os quartos intermediários ( 2.º e 3.º quartos ) apresentaram 38,1 \% e 36,0 \%, respectivamente.
O Gráfico 5.11 traz outros resultados em relação à mesma questão retratada no gráfico anterior, porém por organização acadêmica das IES.
De acordo com os dados do Gráfico 5.11, considerando -se a extensão da prova, o percentual de estudantes que avaliaram a prova como “ Longa ” ou “ Muito longa ”, em relação ao tempo total destinado à sua resolução, foi maior entre os estudantes de universidades, alcançando 38,7 \%, enquanto entre os estudantes de CEFET/IF, a organização acadêmica com a menor incidência, esse índice atingiu 32 \%.
Entre os participantes de centros universitários, o índice atingiu 36,8 \%.
Entre os estudantes de faculdades, 32,8 \%.
O Gráfico 5.12 apresenta os dados por categoria administrativa das IES.
Os dados do Gráfico 5.12 apresentam que o percentual de estudantes que consideraram a prova “ Longa ” ou “ Muito longa ”, em relação ao tempo total destinado à sua resolução, foi maior entre os estudantes de cursos de IES privadas ( 37,6 \% ) do que entre os estudantes de IES públicas ( 34,4 \% ). 5.3 COMPREENSÃO DOS ENUNCIADOS DAS QUESTÕES Organizados em Componente de Formação Geral e em Componente Específico, os dados apresentados nesta seção dizem respeito à forma como os estudantes compreenderam os enunciados das questões da prova, considerando as formulações expressas nesses mesmos enunciados. 5.3.1 COMPONENTE DE FORMAÇÃO GERAL O Gráfico 5.13 apresenta as respostas obtidas para a questão 4 “ Os enunciados das questões da prova na parte de Formação Geral estavam claros e objetivos? ”.
Os resultados apresentados no Gráfico 5.13 apresentam que, na análise regional, a \textbf{percentagem} de estudantes que avaliaram que todos ou a maioria dos enunciados das questões do Componente de Formação Geral estavam “ claros e objetivos ” variou de 79 \%, na região Norte, maior incidência, a 74,8 \%, na região Sudeste, a menor incidência.
O Gráfico 5.14 traz mais resultados sobre a questão tratada no gráfico anterior, porém por quartos de desempenho.
De acordo com os dados do Gráfico 5.14, segundo o desempenho, observa -se que a proporção dos que emitiram essa opinião é crescente, ou seja, \textbf{cresce} conforme o desempenho aumenta.
No quarto superior, a clareza e a objetividade de todos ou da maioria dos enunciados das questões foram percebidas por 86,4 \% dos estudantes.
Já no primeiro quarto inferior, 66,0 \%.
O Gráfico 5.15 traz outros resultados também em relação à mesma questão retratada no gráfico anterior, porém por organização acadêmica das IES.
Os dados do Gráfico 5.15 apresentam que a proporção dos presentes que consideraram que todos ou a maioria dos enunciados das questões do Componente de Formação Geral estavam “ claros e objetivos ” foi maior entre os estudantes de CEFET/IF, chegando a 80,8 \%, enquanto entre os de faculdades, a organização acadêmica com a menor incidência, esse índice atingiu 69,2 \%.
Entre os estudantes de centros universitários, o índice atingiu 75,7 \%.
O Gráfico 5.16 apresenta os dados por categoria administrativa das IES.
Os dados do Gráfico 5.16 apresentam que o percentual de estudantes que consideraram que todos ou a maioria dos enunciados das questões do Componente de Formação Geral estavam “ claros e objetivos ” foi maior entre os estudantes de cursos de IES públicas ( 79,6 \% ) quando comparado aos estudantes de IES privadas ( 75,9 \% ). 5.3.2 COMPONENTE ESPECÍFICO O Gráfico 5.17 apresenta as respostas obtidas para a questão 5 “ Os enunciados das questões da prova na parte de Componente Específico estavam claros e objetivos? ”.
A maioria dos estudantes de todas as grandes regiões brasileiras considerou “ claros e objetivos ” todos ou a maioria dos enunciados das questões do Componente Específico da prova, percentual sempre maior ou igual a 73,5 \%.
Todas as IES apresentando percentuais próximos uns dos outros.
O Gráfico 5.18 traz mais resultados sobre a questão tratada no gráfico anterior, porém por quartos de desempenho.
Os dados do Gráfico 5.18 apresentam que a proporção de estudantes que consideraram os enunciados das questões “ claros e objetivos ” apresenta uma tendência crescente em relação ao aumento de desempenho : mais elevada no quarto superior ( 84,8 \% ), se comparada ao quarto inferior de desempenho ( 64,0 \% ).
O Gráfico 5.19 traz outros resultados também em relação à mesma questão retratada no gráfico anterior, porém por organização acadêmica das IES.
Os dados do Gráfico 5.19 apresentam que a proporção dos presentes que consideraram que todos ou a maioria dos enunciados das questões do Componente Específico estavam “ claros e objetivos ” foi maior entre os estudantes de centros universitários, atingindo 75 \%.
Entre os estudantes de faculdades, a organização acadêmica com a menor incidência, essa proporção atingiu 71,3 \%.
As universidades e CEFET/IF ficaram nos níveis intermediários com 74 \% e 73,5 \%, respectivamente.
O Gráfico 5.20 apresenta os dados por categoria administrativa das IES.
Os dados do Gráfico 5.20 apresentam que o percentual de estudantes que consideraram que todos ou a maioria dos enunciados das questões do Componente Específico estavam “ claros e objetivos ” foi maior entre os estudantes de cursos de IES públicas ( 74,8 \% ) que entre os de IES privadas ( 74,1 \% ). 5.4 SUFICIÊNCIA DAS INFORMAÇÕES/INSTRUÇÕES FORNECIDAS O Gráfico 5.21 apresenta as respostas obtidas para a questão 6 “ As informações/instruções fornecidas para a resolução das questões foram suficientes para resolvê -las? ”.
De acordo com os dados do Gráfico 5.21, e quanto à distribuição de respondentes pelas grandes regiões, observa -se que a proporção de estudantes que consideraram as informações/instruções fornecidas até excessivas ou suficientes em todas ou na maioria das questões foi sempre igual ou superior a 79,3 \%, chegando a 85,2 \%, na região Norte.
O Gráfico 5.22 traz resultados sobre a questão tratada no gráfico anterior, porém por quartos de desempenho.
De acordo com os dados do Gráfico 5.22, levando -se em conta o desempenho dos participantes, observa -se uma tendência crescente das proporções de participantes que avaliaram as informações/instruções como até excessivas ou suficientes em todas ou na maioria das questões.
O percentual foi mais elevado no quarto superior ( 89,8 \% ).
Já no quarto inferior, a suficiência das informações/instruções declarada como até excessiva em todas ou na maioria das questões foi percebida por 72,5 \% dos respondentes.
O Gráfico 5.23 traz outros resultados também em relação à mesma questão retratada no gráfico anterior, porém por organização acadêmica das IES.
De acordo com os dados do Gráfico 5.23, considerando -se as informações/instruções fornecidas para a resolução das questões, a proporção que considerou que essas eram até excessivas ou suficientes em todas ou na maioria das questões foi maior entre os estudantes de CEFET/IF, \textbf{totalizando} 83 \%.
As faculdades foram a organização acadêmica em que houve a menor proporção : 79,1 \%.
O Gráfico 5.24 apresenta os dados por categoria administrativa das IES.
De acordo com os dados do Gráfico 5.24, o percentual de estudantes que consideraram que as informações/instruções fornecidas para a resolução das questões eram até excessivas ou suficientes em todas ou na maioria das questões foi maior entre os estudantes de cursos de IES públicas, 81,8 \%, quando comparado com o percentual das IES privadas : 80,5 \%. 5.5 DIFICULDADE ENCONTRADA AO RESPONDER À PROVA O Gráfico 5.25 apresenta as respostas obtidas para a questão 7 “ Você se deparou com alguma dificuldade ao responder à prova?
Qual? ”.
Nos Gráficos 5.25 a 5.28, apresentados de forma \textbf{sequenciada} mais adiante, são expostos os percentuais de estudantes que apontaram o desconhecimento do conteúdo como “ dificuldade ao responder à prova ”.
O Gráfico 5.25 apresenta as respostas obtidas sobre o desconhecimento do conteúdo como a principal “ dificuldade ao responder à prova ”.
De acordo com os dados o Gráfico 5.25, o percentual de inscritos e presentes que apontaram o desconhecimento do conteúdo como “ dificuldade ao responder à prova ” variou entre de 13,8 \%, na região Nordeste, até 16,4 \%, na região Centro-Oeste.
O Gráfico 5.26 traz resultados sobre o ponto tratado no gráfico anterior, porém por quartos de desempenho.
De acordo com os dados do Gráfico 5.26, em relação aos quartos de desempenho, o desconhecimento do conteúdo foi a opção escolhida por 15,0 \% dos estudantes do segundo quarto inferior e por 14,0 \% do terceiro quarto superior.
A maior porcentagem foi de 16,5 \% no primeiro quarto.
O Gráfico 5.27 traz outros resultados também em relação à mesma questão retratada no gráfico anterior, porém por organização acadêmica das IES.
De acordo com os dados do Gráfico 5.27, na análise por organização acadêmica, o percentual de inscritos e presentes que apontaram o desconhecimento do conteúdo como “ dificuldade ao responder à prova ” é maior nas faculdades ( 17,4 \% ) e, do lado oposto, tem menor incidência nas universidades, com 14,6 \%.
O Gráfico 5.28 apresenta os dados por categoria administrativa das IES.
De acordo com os dados do Gráfico 5.28, em relação às categorias administrativas, o desconhecimento do conteúdo foi a opção escolhida por 16,1 \% dos estudantes, nas IES públicas e 14,8 \% nas IES privadas. 5.6 CONTEÚDOS DAS QUESTÕES OBJETIVAS DA PROVA O Gráfico 5.29 apresenta as respostas obtidas para a questão 8 “ Considerando apenas as questões objetivas da prova, você percebeu que não estudou ainda a maioria desses conteúdos? ”.
Na análise por grande região, a proporção de respondentes que escolheram a opção a qual indica que não estudaram ainda a maioria desses conteúdos foi pequena.
Observa -se que, nas regiões Sul ( 6,6 \% ) e Norte ( 7,4 \% ) foram menores.
Nas regiões Sudeste e Nordeste, a situação foi \textbf{inversa}, com proporções mais altas, correspondendo a, respectivamente, 8,5 \% e 8,2 \%.
A região Centro-Oeste ficou com a proporção de 8 \%.
O Gráfico 5.30 traz resultados sobre a questão tratada no gráfico anterior, porém por quartos de desempenho.
De acordo com os dados do Gráfico 5.30, considerando -se separadamente as opiniões de estudantes dos quatro quartos de desempenho, observa -se uma tendência decrescente, no quarto inferior, 14,2 \% ofereceram como resposta que não estudaram ainda a maioria desses conteúdos, enquanto 3 \% do quarto superior informaram a mesma resposta.
O Gráfico 5.31 traz outros resultados também em relação à mesma questão retratada no gráfico anterior, porém por organização acadêmica das IES.
De acordo com os dados do Gráfico 5.31, na análise por organização acadêmica, a proporção de respondentes que escolheram a opção a qual indica que não estudaram ainda a maioria desses conteúdos foi pequena.
Observa -se que as proporções nas faculdades ( 11,1 \% ) e nos centros universitários ( 7,9 \% ) foram maiores.
Nas universidades e nos CEFET/IF, a situação foi \textbf{inversa} : a proporção foi menor 7,8 \% e 6 \%, respectivamente.
O Gráfico 5.32 apresenta os dados por categoria administrativa das IES.
De acordo com os dados do Gráfico 5.32, em relação às categorias administrativas, observa -se que, nas IES públicas, 5,3 \% ofereceram como resposta que “ não estudaram ainda a maioria desses conteúdos ”.
Enquanto os estudantes de IES privadas apresentaram maior desconhecimento quanto aos conteúdos ( 8,1 \% ). 5.7 TEMPO GASTO PARA CONCLUIR A PROVA O Gráfico 5.33 apresenta as respostas obtidas para a questão 9 “ Qual foi o tempo gasto por você para concluir a prova? ”.
Considerando -se as cinco grandes regiões brasileiras, os percentuais dos estudantes que utilizaram entre duas e quatro horas para finalizar a prova foram, nas regiões Norte e Nordeste, respectivamente, 76,4 \% e 76,5 \%.
Nas regiões Sul e Centro-Oeste, os percentuais de estudantes que \textbf{dispensaram} entre duas e quatro horas para concluir a prova foram menores, 72,5 \% e 70,0 \%, respectivamente.
A região Sudeste ficou com o percentual de 74,5 \%.
O Gráfico 5.34 traz resultados sobre a questão tratada no gráfico anterior, porém por quartos de desempenho.
De acordo com os dados do Gráfico 5.34, uma vez analisadas as alternativas escolhidas pelos estudantes que se situaram nos diferentes quartos de desempenho, observa -se uma tendência crescente : uma proporção maior de participantes no quarto superior ( 81,7 \% ) declarou ter gastado entre duas e quatro horas para concluir a prova em comparação com a proporção de participantes nos quartos inferiores, atingindo 62,7 \% no primeiro quarto inferior.
O Gráfico 5.35 traz outros resultados também em relação à mesma questão retratada no gráfico anterior, porém por organização acadêmica das IES.
De acordo com os dados do Gráfico 5.35, na análise por organização acadêmica, os percentuais dos que utilizaram entre duas e quatro horas para finalizar a prova foram de 74,1 \% nas universidades ; 76,8 \% no CEFET/IF e 73 \% nos centros universitários.
O menor percentual foi entre os estudantes de faculdades, com 71,8 \% dos alunos necessitando de duas a quatro horas para concluir a prova.
O Gráfico 5.36 apresenta os dados por categoria administrativa das IES.
De acordo com os dados do Gráfico 5.36, em relação às categorias administrativas, observa -se que um percentual maior de estudantes nas IES públicas ( 76,8 \% ) declarou ter gasto entre duas e quatro horas para concluir a prova, enquanto nas IES privadas esse percentual foi de 73,5 \%.
Os resultados do Enade 2022 apresentados neste capítulo contribuem para o entendimento de que a consideração da percepção dos estudantes concluintes sobre a prova do Enade é importante tanto para a compreensão do seu desempenho como para a reflexão sistemática e constante a respeito de determinadas realidades relacionadas a esse mesmo desempenho no país.
Por outro lado, como todos os elementos do trabalho \textbf{pedagógico} institucional apresentam -se diretamente relacionados com o desempenho dos estudantes, também ganha relevância conhecer e analisar os resultados da percepção dos \textbf{coordenadores} de curso sobre a mesma prova.
As informações expressas por estudantes e \textbf{coordenadores} sobre questões diversas associadas aos dois componentes da prova – Formação Geral e Componente Específico – podem agregar valor em ações institucionais voltadas para o aperfeiçoamento do instrumento \textbf{avaliativo}, considerando os vários elementos envolvidos em sua elaboração.
Neste contexto, é relevante que os resultados das percepções de estudantes e \textbf{coordenadores} sobre a prova do Enade ganhem peso como um elemento a mais nas análises da oferta dos cursos nas modalidades presencial e a distância, em nível nacional.
Espera -se, assim, que, em uma visão formativa de avaliação, o conhecimento e a análise desses resultados possam contribuir para a construção de alternativas institucionais que venham a induzir a qualidade da educação superior.

% matched lemmas: arredondamento, avaliativo, coordenador, crescer, disjunto, dispensar, estimador, interseção, inverso, numérico, ordenar, parâmetro, pedagógico, percentagem, questionário, recorte, reiterar, reprodução, sequenciado, subconjunto, totalizar, vidar
\end{textsample}
